\documentclass[11pt,a4paper]{article}
\usepackage{fontspec}
\setmainfont{Noto Serif CJK KR}
\setsansfont{Noto Sans CJK KR}
\usepackage{amsmath,amssymb,bm}
\usepackage[margin=2.3cm]{geometry}
\usepackage{graphicx,booktabs,array}
\usepackage[hidelinks]{hyperref}
\usepackage{caption}
\captionsetup{font=small,labelfont=bf}
\setlength{\parskip}{3pt}
\newcommand{\ehat}{\hat{e}}
\newcommand{\Ghat}{\hat{G}}
\newcommand{\Jhat}{\hat{J}}

\title{\textbf{절단 없는 Bianchi 배경 Boltzmann--Einstein 솔버}\\[3pt]
\large 편광 포함 비섭동 격자 운동론 --- 결과 보고서}
\author{\normalsize \texttt{bianchi/q} 85차 코드 기준 (자동 생성 수치)}
\date{}

\begin{document}
\maketitle

\begin{abstract}\noindent
11개 Bianchi 유형 전부에 대해 균질·비섭동 배경에서 Einstein 방정식과 광자
Boltzmann 방정식을 \textbf{모멘트 절단 없이} 푸는 솔버의 결과를 보고한다.
상태는 위상공간 격자 위의 분포함수 자체이고, 유체 닫힘·$\ell_{\max}$ 절단·
근사 절환(tight-coupling $\leftrightarrow$ 완전결합)이 모두 없다.
이번 판에서 새로 닫힌 것은 (i) 광자 \textbf{편광}의 기저-없는 텐서 캐리어와
그 정확 충돌 지수, (ii) 편광과 \textbf{전자 벌크운동(부스트)}의 결합,
(iii) 외부 재결합 이력 $x_e(z)$ 를 적분기까지 잇는 배선이다.
주요 결과: Thomson 충돌 지수가 편광까지 포함해 \textbf{정확히 3항}으로 닫히고
($\mathcal K$ 가 rank 9, 고유값 $(1,\,7/10,\,1/2)$) 강성이 소멸하며;
편광/무편광 차가 $\ell=2$ 만 통해 들어오는 \textbf{닫힌형}으로 주어지고;
부스트의 \textbf{정준 수송 스크린 사상}에 추가 회전이 없음 ($\psi=0$) 이
기계정밀로 확인되었다. 편광 채널을 켜도 강도 $I$ 채널은 스칼라 경로를 \textbf{비트
단위로} 재현한다.
\end{abstract}

\begin{center}
\fbox{\parbox{0.95\textwidth}{\small
\textbf{범위 경계 (가장 큰 것).}\;
이온화 이력 $x_e(z)$ 는 \textbf{외부에서 공급}된다.  재결합과 재이온화는 아직
자기일관하게 풀리지 않는다.  현재 어댑터는 외부 $x_e(z)$ 를 매 스텝 소비하지만,
그 $x_e(z)$ 가 FLRW 팽창과 등방 복사장 아래에서 만들어진 것이라면 이는 엄밀히
\textbf{external-history lane} 이다.  자체 재결합은 $H_{\rm Bianchi}$,
$f_\gamma(E,\ehat)$, 원자 준위, Ly$\alpha$ 복사수송이 함께 되먹임되어야 하며 ---
이것은 이 솔버의 마지막 5\,\% 가 아니라 사실상 \textbf{별도의 physics engine} 이다.
}}
\end{center}

\section{솔버의 구조 --- 벗은 것과 남은 것}

시간변수는 $d\tau = H\,dt$, 상태는 Wainwright--Ellis 팽창정규화 변수
$(\Sigma_{ab}, N_{ab}, A_a, \ln H)$ 와 위상공간 격자 위의 분포다.
\textbf{벗은 것}(근사 부재)과 \textbf{남은 것}(이산화)을 계약
(\texttt{bianchi/q/contract.py})에 기계가독으로 못박고, 항목마다 그것을 재는
시험 ID를 붙였다.

\begin{center}\small
\begin{tabular}{@{}ll@{}}
\toprule
\textbf{벗은 것} & \textbf{무엇이 그것을 보장하는가} \\
\midrule
모멘트 절단 ($\ell_{\max}$, 계층 닫힘) & 상태가 $f$ 자체 \\
충돌 선형화 (소-$v$ 전개) & 정확 광행차; $O(v^2)$ 지수 실측 \\
섭동 전개 & $O(1)$ 이방성에서 직접 적분 \\
근사 절환 & 정확 지수 $\Rightarrow$ \textbf{충돌 부분스텝}의 안정성 강성 제거 \\
유체 닫힘 & $\rho,q_a,\pi_{ab}$ 는 격자 구적의 \emph{출력} \\
\bottomrule
\end{tabular}
\end{center}

\medskip\noindent
\textbf{공변 운동량 프레임.} $p = M(\tau)q$, $dM/d\tau = -(I+\Sigma-\varepsilon R)M$
로 두면 팽창·전단·회전이 \emph{정확히} 흡수되어 방향격자 보간이 0회가 된다.
곡률 $(N,A)$ 가 남기는 방향공간 흐름만 반-라그랑주로 처리한다.
상태를 $\ln\Ghat$ 로 들고 모멘트를 $(\ln\rho,\,q/\rho,\,\pi/\rho)$ 로 돌려주므로
깊은 붕괴에서 $\rho$ 가 배정도 범위를 넘어도 벽이 없다.

\section{정확 충돌 지수와 강성 소멸}

Thomson 핵은 정확히 $\ell\le2$ 대역제한이므로 ($k_\ell=(1,0,1/10,0,\dots)$)
스칼라 충돌 지수가 3항으로 닫힌다:
\begin{equation}
e^{xC}f = P_0 f + e^{-x}\bigl(f-P_0f-P_2f\bigr) + e^{-0.9x}P_2 f,
\qquad x=\nu\Delta\tau .
\end{equation}
이것이 근사 절환이 필요 없는 이유다: $x\to\infty$ 가 안정할 뿐 아니라 정확히
$P_0 f$ 로 떨어진다 (그림~\ref{fig:stiff}a). Sinkhorn·행렬지수 기계가 전부 불필요하다.

\paragraph{주장의 범위.} 정확히 말하면 제거된 것은 \textbf{충돌 부분스텝의 안정성
강성}이다 (\emph{stability stiffness eliminated in the collision step}).  수송과
충돌은 비가환이고 $\nu_\tau$ 는 재결합에서 급격히 변하므로 \textbf{분할 정확도
제약은 남는다} --- 오차예산에서 Strang 이 정확히 2차로 나오는 것이 바로 그
증거다.  즉 ``$\Delta\tau$ 를 안정성 때문에 줄일 필요가 없다'' 는 참이고,
``정확도 때문에도 줄일 필요가 없다'' 는 거짓이다.

\section{편광 --- 기저-없는 캐리어와 rank-9 정확 지수}

캐리어는 기저-없는 에르미트 3-텐서 $J_{ab}(\ehat)$, 제약 $J_{ab}\ehat^b=0$.
실수 대칭 6 (강도 $I$ + 선형편광) + 반대칭 3 (원편광 $V$) = \textbf{9 성분}.
\emph{Stokes 부호 규약이 필요 없다} --- 기저를 쓰지 않기 때문이다 (IAU/CMB 의
$U$ 부호, $V$ 부호, 스크린 방향 문제가 아예 발생하지 않는다).

편광 Thomson 연산자
\begin{equation}
(KJ)_{ab}(\ehat') = \frac{3}{8\pi}\int d\Omega\;
\Pi_{ac}(\ehat')\,J_{cd}(\ehat)\,\Pi_{db}(\ehat'),\qquad \Pi=\delta-\ehat\ehat
\end{equation}
에서 투영자가 $\ehat'$ 에만 의존하므로 $\ehat$ 적분이 $M_{cd}=\int J_{cd}d\Omega$
(9 성분) 로 축약된다. 즉 $K$ 는 \textbf{정확히 rank 9} 이고 $\ell$-분해가 필요 없다.
구면 항등으로
\begin{equation}
\mathcal K[M]=\tfrac{7}{10}M+\tfrac{1}{10}\mathrm{tr}(M)\delta\ \ (\text{대칭}),
\qquad \mathcal K[A]=\tfrac12 A\ \ (\text{반대칭}),
\end{equation}
고유값이 $(1,\,7/10,\,1/2)$ --- 각각 광자수 보존, 무대각 대칭 5차원, $V$ 3차원 ---
이고 정확 지수가 다시 \textbf{3항}이 된다. 이 고유값은 하드코딩하지 않고 매번
구적으로 재계산해 대조한다 (그림~\ref{fig:stiff}b).

\begin{figure}[htbp]\centering
\includegraphics[width=\textwidth]{fig/f4_stiff.pdf}
\caption{(a) 강성 소멸: 정확 3항이 $x\to\infty$ 에서 고유값-1 부분공간으로
정확히 사영된다 (고정점 잔차 $2.2\times10^{-16}$). 접근율은 stf 간극 $1-0.7$ 이 정하는
$e^{-0.3x}$. (b) $\mathcal K$ 고유값을 각 해상도별로 구적 재계산한 오차.
$n_\theta$ 를 올리면 격자 가중치의 반올림 때문에 \emph{오차가 커진다} ---
``해상도를 올리면 좋아진다''는 직관이 이 게이트에는 거짓이다.}
\label{fig:stiff}
\end{figure}

\section{결과 1 --- 편광/무편광 차의 닫힌형 ($E\!\to\!I$ 되먹임)}

편광 연산자의 트레이스를 스칼라 3항과 비교하면 차가 남는다. 이전 판에서는
그것이 구적 앨리어싱인지 물리인지 판정하지 못해 미해결로 두었다. 이번에 닫혔다:
\begin{equation}
\boxed{\ \Delta I(\ehat)=\Bigl[\tfrac67 e^{-x}+\tfrac17 e^{-0.3x}-e^{-0.9x}\Bigr]
(P_2 I_0)(\ehat),\qquad x=\nu\Delta\tau\ }
\end{equation}
$\ell=2$ 블록이 $\bigl(\begin{smallmatrix}1/10 & -\sqrt6/10\\ -\sqrt6/10 & 3/5\end{smallmatrix}\bigr)$
(trace $7/10$, det $0$ $\Rightarrow$ 고유값 $\{0,7/10\}$) 라는 데서 나온다.
소 $x$ 에서 $\Delta I\simeq 0.03\,x^2 (P_2I_0)$.

\textbf{왜 결정적인가.} 이 식의 모든 단계가 구적 가중치에 대해 \emph{점별 대수}이므로
임의의 구적에서 정확하다 --- 결손 $2\times4$ 격자에서도 $2.4\times10^{-16}$.
즉 구적 오차가 전혀 없다. 앨리어싱 가설의 결정적 기각이다. 차는 $\ell=2$ 를
통해서만 들어오고 ($\ell=0,1,3,4$ 는 $\lesssim10^{-14}$), $\ell=2$ 에서만
$1.73\times10^{-2}$ 로 유의미하다 (그림~\ref{fig:fb}).

문헌의 $k_2=1/10$ 은 이 블록의 \emph{강도-강도} 성분이고 $7/10$ 은 그 trace 다 ---
모순이 아니라 같은 행렬의 두 얼굴이다. 준정적 소거로 $3/4$ 가 나와
Kaiser/Hu 규약과 일치한다.

\begin{figure}[htbp]\centering
\includegraphics[width=\textwidth]{fig/f3_feedback.pdf}
\caption{(a) 세 가지 각 해상도에서 잰 $\max|\Delta I|$ 와 닫힌형(검은 선).
격자가 달라도 같은 곡선 위에 놓인다 --- 물리이지 앨리어싱이 아니다.
(b) 입력 다중극별 편광/무편광 차: $\ell=2$ 만 유의미하다.}
\label{fig:fb}
\end{figure}

\section{결과 2 --- 정준 수송 스크린 사상에서 $\psi=0$ (부스트)}

전자 벌크속도 $v_b$ 로의 부스트에서 스크린 대표원 $W$ ($W^0=0$, $W\!\cdot\!p=0$) 를
4-벡터로 나른 뒤 게이지를 복원하면
$W'_a=(\Lambda W)_a-\bigl((\Lambda W)^0/E'\bigr)p'_a$ 다. 이 사상
$S:\text{screen}(\ehat)\to\text{screen}(\ehat')$ 을 산란면 적응기저
($m_1\propto P_\perp[v]$, $m_2=\ehat\times m_1$) 에서 SVD 한 결과:

\begin{center}\small
\begin{tabular}{@{}lr@{}}
\toprule
$|\text{배율}-1|$ & $4.4\times10^{-16}$ \\
$|\psi|$ (추가 회전) & $2.2\times10^{-15}$ \\
왕복(정지계 왕복) 항등 & $\lesssim10^{-15}$ \\
무편광 $\to$ 무편광 & $\lesssim10^{-15}$ \\
\bottomrule
\end{tabular}
\end{center}

\noindent 즉 \textbf{부스트는 추가 회전을 만들지 않고 배율은 정확히 1} 이다.
계획 단계의 예상(``광행차가 스크린을 돌린다'')은 \emph{기각}되었다. 광행차는
$\ehat$ 를 돌리지만 스크린 \emph{안에서의} 회전은 만들지 않는다.

배율이 1 인 이유는 해석적으로 명확하다: $\Lambda$ 는 계량을 보존하고 게이지
복원에서 빼는 것이 \emph{널} 벡터 $p'$ 의 배수이므로
$|W'-\alpha p'|^2=|W'|^2-2\alpha(W'\!\cdot\!p')+\alpha^2(p'\!\cdot\!p')$ 에서
마지막 항은 $p'$ 이 널이라 0, 가운데 항은 $W\!\cdot\!p=0$ 의 부스트 불변성으로 0 이다.
``배율이 $1/D$ 일 것''이라는 순진한 예상은 틀렸다 --- 도플러 $D$ 는
\emph{진폭}($E'=DE$)에 붙지 스크린 벡터에 붙지 않는다 (그림~\ref{fig:boost}).

\paragraph{주장의 범위.} 여기서 $\psi=0$ 은 \textbf{정준 수송 스크린 사상}에 대한
진술이다 --- 스크린 대표원 $W$ 를 Lorentz 수송한 뒤 게이지를 복원해 얻은 사상 $S$ 를,
산란면 적응기저 ($m_1\propto P_\perp[v]$) 에서 읽은 값이다.  \textbf{임의의 외부
tetrad 나 전역 편광 기저를 고정하면 Wigner phase 가 다시 나타날 수 있다} --- 그것은
기저 선택의 성질이지 $S$ 의 성질이 아니다.  이 코드가 기저-없는 텐서 캐리어를 쓰기
때문에 실제 계산에서 그 phase 가 등장하지 않는 것이고, ``부스트가 편광면을 절대
돌리지 않는다'' 는 더 넓은 명제를 주장하는 것이 아니다.  논문 표기는
\emph{``$\psi=0$ for the canonical transported screen map''} 로 좁힌다.

\textbf{귀결.} 편광 $\times$ 부스트는 새 기계가 필요 없다.
$J'_{ab}(\ehat')=S_a{}^cS_b{}^dJ_{cd}(\ehat)$ 이고 구현은
\emph{정지계로 부스트 $\to$ rank-9 3항 $\to$ 되부스트} 세 단계다.

\begin{figure}[htbp]\centering
\includegraphics[width=\textwidth]{fig/f5_boost.pdf}
\caption{(a) 무작위 600 조합 $(v_b,\ehat)$, $|v_b|\le0.85$ 에서 잰 스크린 사상의
회전각·배율·이방도. 전부 기계정밀 바닥에 붙어 있다. (b) 도플러 인자는 진폭에만
작용한다 --- 스크린 사상의 배율은 $|v_b|$ 와 무관하게 1.}
\label{fig:boost}
\end{figure}

\section{결과 3 --- $O(1)$ 이방 배경이 만드는 편광 (섭동전개 없이)}

선형 섭동론을 쓰지 않고, $\Sigma$ 가 $O(1)$ 인 배경에서 Thomson 산란이 만드는
편광을 직접 계산한다. 등방 입사는 편광을 정확히 0 만들고 ($\lesssim10^{-14}$),
사원극 이방성이 있을 때만 편광이 생긴다 --- 물리 서명이 격자 위에서 그대로
재현된다. 그림~\ref{fig:polmap} 은 $\Sigma=\mathrm{diag}(0.35,-0.15,-0.20)$,
$N=\mathrm{diag}(0.4,-0.3,0.1)$ 배경에서 $\nu\Delta\tau=1$ 일 때의 편광도 분포와
광학깊이에 대한 응답을 보인다 (최대 편광도 $0.034$).

정량 관측량으로 쓰는 것은 기저-없는 편광도 $p=\sqrt{Q^2+U^2+V^2}/I$ 뿐이다.
Stokes 다중극과 $E/B$ 분해는 \texttt{\_diagnostic} 접미사를 달아 \emph{정성 진단
전용}으로 격리했다.  \textbf{격리하는 이유는 구현 쪽이다.}  $E/B$ 분해 \emph{자체는}
선형 섭동론 전용 언어가 아니다 --- spin-2 장의 전역 분해는 비섭동적으로도 정의된다.
문제는 현재 쓰는 국소 스크린 게이지 ($\arg\min|\ehat|$ 기저) 가 구면 위에서
불연속이라, 그 위에서 정의한 $Q/U$ 가 각 해상도로 \textbf{수렴하지 않고}
(실측 $-4.8\times10^{-4}\ldots-2.3\times10^{-4}$; $I_2$ 는 9자리 안정) 따라서
전역 spin-2 분해가 \textbf{정량적으로 안정적이지 않다}는 데 있다.  연속적인 스크린
게이지 (또는 spin-weighted 구면조화 기반) 를 도입하면 정량화할 수 있다 --- 그때까지만
diagnostic 이다.

\begin{figure}[htbp]\centering
\includegraphics[width=\textwidth]{fig/f6_polmap.pdf}
\caption{(a) $O(1)$ 이방 배경에서 한 번의 Thomson 광학깊이가 만든 편광도 분포.
(b) 편광 생성은 저광학깊이에서 $\nu\Delta\tau$ 에 선형이고 $\nu\Delta\tau\gtrsim1$ 에서 포화한다 --- 등방화가 사원극을 지우기 때문이다.}
\label{fig:polmap}
\end{figure}

\section{결과 4 --- 편광 채널이 강도 채널을 비트 단위로 보존한다}

편광을 켜도 기존 스칼라 솔버가 깨지지 않아야 한다. 캐리어를
\emph{진폭 $\times$ 모양}으로 쪼개면 구조적으로 보장된다:
\begin{equation}
J_{ab}(\ehat)=\Ghat(\ehat)\,\Jhat_{ab}(\ehat),\qquad
\mathrm{tr}\,\Jhat=1,\quad \Jhat_{ab}\ehat^b=0 .
\end{equation}
진폭 $\ln\Ghat$ 는 기존 경로 그대로 두고 모양 $\Jhat$ 만 새로 나른다.
지표 수송이 $\mathrm{tr}(dJ)=0$ 을 정확히 만족하므로 진폭을 건드리지 않고,
충돌은 $J$ 에 선형이라 공통 인수를 밖으로 뺄 수 있다 (로그공간 범위 안전).

실측: Bianchi I, V, VIII, IX 전부에서 자유흐름의 $I$ 채널이 스칼라 경로와
\textbf{정확히 0.0} 만큼 다르다. 무편광 상태는 곡률이 있어도 무편광으로 남고,
그 잔차 $5.5\times10^{-09}$ 는 \emph{각 해상도와 무관}하며 $\Delta\tau$ 로만 3차 수렴한다.

\begin{figure}[htbp]\centering
\includegraphics[width=\textwidth]{fig/f7_p8.pdf}
\caption{(a) 편광 run 의 $I$ 채널이 스칼라 run 을 비트 단위로 재현한다
(4개 유형 전부 정확히 0). (b) 자유흐름이 만드는 허위 편광. 채택된 해석적
분리(파랑)는 각 해상도에 무관하고, 기각된 순진한 remap(빨강)은 각 축으로만
$\sim1.7$ 차 수렴했다 --- $\Delta\tau$ 를 줄여도 안 줄어드는 ``누락된 항''의 서명.}
\label{fig:p8}
\end{figure}

\section{결과 5 --- 외부 재결합 이력의 실제 소비}

이전 판까지 이온화 이력 $x_e(z)$ 는 객체 수준에서 \emph{접합}만 되어 있었고
적분기에는 상수 $\nu$ 가 들어갔다. 이번에 그 구간을 닫았다. $\tau$ 의 정의에서
$d\ln a/d\tau=1$ 이므로
\begin{equation}
1+z(\tau)=(1+z_0)\,e^{-(\tau-\tau_0)}
\end{equation}
가 \emph{정확}하고 (이방 배경에서도 $a$ 는 $H=\dot a/a$ 로 정의된 평균
척도인자이므로 근사가 아니다), 커널이 매 스텝 $\nu_\tau=\sigma_T n_e c/H$ 를
계산한다. Python 스텝 루프는 여전히 0회다.

한 run 안에서 $\nu_\tau$ 가 $5.6\times10^{-14}$ 에서 $356.1$ 까지 \textbf{약 16
자릿수} 변한다 --- 상수 $\nu$ 로 근사하면 통째로 버려지는 정보다. 배선 전체는
물리 항등 $d\tau_{\rm opt}=\nu_\tau d\tau$ 로 검증된다: 솔버가 적분한
$\int\nu_\tau d\tau=170.97584$ 가 외부 모듈의 광학깊이 $170.97584$ 와 상대
$2.1\times10^{-09}$ 로 일치한다 (그림~\ref{fig:rate}).

\begin{figure}[htbp]\centering
\includegraphics[width=\textwidth]{fig/f1_rate.pdf}
\caption{(a) 적분기에 실제로 들어가는 $\nu_\tau(z)$ 와 $x_e(z)$ (점선).
(b) 가시함수와, 솔버가 적분한 광학깊이 대 외부 모듈 값의 대조.}
\label{fig:rate}
\end{figure}

\subsection*{$H$ 규약의 정직성}

$\nu_\tau$ 의 분모 $H$ 를 무엇으로 잡느냐는 이방 배경에서 물리적 의미가 있다.
기본값은 \emph{모델 자체의} $H$ 이지만 차원 앵커가 $\Lambda$CDM $H(z_0)$ 이므로
$\tau_0$ 에서 두 규약의 $\nu$ 는 정확히 같다. 즉 모델 규약은 $\tau_0$
\emph{이후의 모양}만 자기일관하게 준다. 이 구간에서 두 규약은 최대
$33.8\%$ 차이가 난다 (그림~\ref{fig:hist}b). 자기일관 앵커는 모델 자신의
Gauss 면 $\Omega_0$ 에서 유도하며 ($\Lambda$CDM 앵커의 $0.389$ 배) --- 무시할
크기가 아니다. 이 사실을 계약과 \texttt{summary()} 출력에 명시해 숨은 가정이
되지 않게 했다.

\begin{figure}[htbp]\centering
\includegraphics[width=\textwidth]{fig/f2_history.pdf}
\caption{(a) $x_e$ 를 약 10\% 낮춘 이력이 결과를 바꾼다 (상대 $7.1\times10^{-05}$) ---
이력의 \emph{내용}이 소비된다는 증거. 상수 $\nu$ 배선에서는 두 곡선이 완전히
겹친다. (b) 이방 배경에서 모델 $H$ 와 $\Lambda$CDM $H$ 가 주는 $\nu_\tau$ 의 비.}
\label{fig:hist}
\end{figure}

\section{성능 --- Mode B 편광 커널}

운동량 격자를 함께 푸는 Mode B 에서 편광 충돌을 Rust 로 이식했다. Thomson 핵은
에너지 교환이 없으므로 반경 슬라이스마다 같은 연산자이고, 따라서 각 슬라이스
결과가 Mode A 커널과 \emph{정확히} 같아야 한다 (게이트). Rust와 Python 참조가
상대 $3.6\times10^{-16}$ 로 일치하고 평균 $3.4$ 배 빠르다 (그림~\ref{fig:modeb}).

\begin{figure}[htbp]\centering
\includegraphics[width=\textwidth]{fig/f8_modeb.pdf}
\caption{(a) Mode B 편광 충돌 비용. (b) Rust/Python 비트급 파리티.}
\label{fig:modeb}
\end{figure}


\section{11 Bianchi 유형 커버리지 (transport core 동결 기준선)}

초록이 ``11종 전부'' 를 주장하므로 유형별 smoke/pass 표를 기준선으로 고정한다.
12 항목 (11 표준 유형 + 예외형 VI$^*_{-1/9}$) 전부 분류 왕복 $\cdot$ 유한성 $\cdot$
Jacobi 기계정밀 $\cdot$ Codazzi 수렴을 통과한다.  격자 $16\times32$, 200 스텝,
$\tau: 0\to-0.5$, 유형별 튜닝 없음.

\begin{center}\small
\begin{tabular}{@{}lllllll@{}}
\toprule
유형 & class & Gauss & Codazzi & 수렴차수 & Jacobi & 상태 \\
\midrule
I & A & 1.3e-13 & 4.3e-17 & 기계정밀 & 0e+00 & PASS \\
II & A & 4.6e-06 & 2.3e-05 & 3.63 & 0e+00 & PASS \\
IV & B & 3.0e-05 & 1.9e-04 & 2.28 & 2e-19 & PASS \\
V & B & 3.0e-05 & 1.9e-04 & 2.29 & 0e+00 & PASS \\
VI$_{0}$ & A & 8.4e-06 & 2.8e-05 & 3.66 & 0e+00 & PASS \\
VI$_{h}$ & B & 1.3e-05 & 8.0e-05 & 2.50 & 3e-20 & PASS \\
VII$_{0}$ & A & 5.8e-06 & 2.8e-05 & 3.61 & 0e+00 & PASS \\
VII$_{h}$ & B & 1.3e-05 & 8.3e-05 & 2.50 & 3e-20 & PASS \\
VIII & A & 1.2e-05 & 5.4e-05 & 3.52 & 0e+00 & PASS \\
IX & A & 9.8e-07 & 4.2e-06 & 3.72 & 0e+00 & PASS \\
III & B & 2.8e-05 & 1.8e-04 & 2.28 & 1e-19 & PASS \\
VI$^*_{-1/9}$ & B & 1.4e-05 & 1.3e-04 & 2.70 & 3e-19 & PASS \\
\bottomrule
\end{tabular}
\end{center}

\paragraph{★★ 이 표가 내 각주를 반증했다.} 처음엔 ``class A 는 Codazzi 가
기계정밀'' 이라고 적었는데 \textbf{거짓}이다.  기계정밀인 것은 \textbf{Bianchi I
뿐}이고 (곡률항이 아예 없다), 곡률이 있는 class A 도 $10^{-5}$ 수준의 잔차를
갖는다.  차이는 \textbf{크기가 아니라 수렴 차수}에 있다:

\begin{center}\small
\begin{tabular}{@{}lll@{}}
\toprule
묶음 & 유형 & Codazzi 수렴 차수 \\
\midrule
평탄 & I & 기계정밀 ($2\times10^{-16}$) \\
class A 곡률형 & II, VI$_0$, VII$_0$, VIII, IX & \textbf{3.52 -- 3.72} \\
class B & IV, V, VI$_h$, VII$_h$, III, VI$^*$ & \textbf{2.28 -- 2.70} \\
\bottomrule
\end{tabular}
\end{center}

\noindent class A 곡률형이 고차인 것은 T4 (N-항의 $\ell\le1$ 모멘트가 정확히 0) 가
저차 오차를 막기 때문이고, class B 가 저차인 것은 A-항이 $\ell=1$ 에 \textbf{직접}
작용하기 때문이다 --- Q5b 가 등재한 $\mu$-꺾임 축(2.45)과 \textbf{같은 축}이다.
그래서 합격 판정을 절대 임계가 아니라 \textbf{수렴 차수 $\ge1.5$} 로 두었다.
절대 임계를 쓰면 곡률형 전부가 해상도만 낮으면 탈락하고 해상도를 올리면 통과한다 ---
물리가 아니라 격자를 재는 게이트가 된다.

\textbf{범위 주의.} 이 표는 \textbf{구속식과 유한성}을 재는 smoke 매트릭스다.
유형별 물리 결과(각 유형의 Kasner 지수, Mixmaster 진동)를 검증하는 표가 아니다 ---
그 검증은 Tier2 정확해 대조와 I2b 지수 게이트가 담당한다.

\section{검증 체계와 오차예산}

수렴 축의 성격을 세 종류로 분리해 자동 생성 문서에 등재한다.

\begin{center}\small
\begin{tabular}{@{}lll@{}}
\toprule
축 & 성격 & 실측 \\
\midrule
시간 (RK4) & 대수 4차 & 3.99,\;3.99 \\
분할 (Strang) & 대수 2차 & 2.00,\;2.00 \\
반경 (시프트) & 대수 $\ge6$ & 6.5--7.2 \\
각 (GL$\times$균등, 매끄러움) & 스펙트럴 & $4.2\times10^{-14}\to0$ \\
각 (공변 꺾임, 잔여 이류) & \textbf{대수 2.45} & 전체에서 유일한 대수 축 \\
Mode B 각 / 반경 정의역 & 스펙트럴 & 감쇠율 0.60 / 1.74 per node \\
충돌 (정확 3항) & \textbf{정확} & $2.8\times10^{-16}$ \\
편광 지표 수송 (RK4) & 대수 4차 & 3.99 \\
편광 충돌 (rank-9 3항) & \textbf{정확} & $2.8\times10^{-16}$ \\
P8 합성 (무편광 유지) & 대수 3차 & 3.00,\ \emph{각 해상도 무관} \\
\bottomrule
\end{tabular}
\end{center}

\noindent 요점: \textbf{편광을 켜도 새로운 수렴 축이 생기지 않는다.} 지표 수송은
기존 시간축과 같은 적분기를 쓰고, 편광 충돌은 스칼라와 마찬가지로 오차가 아니라
정확이다.

\section{반증 기록 --- 이 코드가 틀렸던 지점들}

통과한 게이트뿐 아니라 \emph{기각된 후보}와 \emph{무력했던 게이트}를 같은
비중으로 기록한다.

\begin{itemize}\itemsep2pt
\item \textbf{고정 tetrad 프레임 기각.} 축비 $2.2\times10^{28}$ 의 Mixmaster
붕괴에서 물리 각분포가 폭 $\sim10^{-28}$ 의 연필빔이 되어 고정격자로는
\emph{표현 자체가 불가능}했다 (보간 200배 감축이 발산을 3자릿수밖에 못 줄였다).
공변 프레임으로 전환해 보간 0회를 달성했다.

\item \textbf{기억한 Koszul 식이 틀렸다.} $\tfrac12(C_{abc}+C_{cab}-C_{bca})$ 는
옳은 식과 정확히 $C_{abc}$ 만큼 다르다. $C_{abc}$ 가 $(a,b)$ 반대칭이 아니라
$\Gamma_{(ab)c}\neq0$ 이 되어 계량 적합이 깨지지만, 동시에 $(b,c)$ 에는 반대칭이라
\emph{측지선에서 통째로 빠진다}. 그래서 운동량 대조는 $2.1\times10^{-14}$ 로
통과하고 편광만 $3.2\times10^{-1}$ 로 틀렸다. $\Rightarrow$ ``운동량이 맞으니
접속이 맞다''는 추론은 성립하지 않는다.

\item \textbf{편광도 함수가 $1/\sqrt2$ 만큼 과소보고.}
$\|J-(I/2)\Pi\|_F=\sqrt{Q^2+U^2+V^2}/\sqrt2$ 이므로 $\sqrt2$ 를 곱해야 한다.
100\% 편광 상태에서 $0.7071$ 을 돌려주고 있었고 임계형 시험만 있어서 아무것도
걸리지 않았다.

\item \textbf{실패할 수 없는 게이트 5건.} 강성 시험이 비정규수 바닥 때문에 두 값을
비트 동일로 비교했고, 텐서율 시험은 구현 몸통을 복사해 정확히 0 을 쟀으며,
rank-9 시험은 ``$\times0$'' 사장코드였다. 전부 판별력 있는 진술로 교체했다.

\item \textbf{시간의존 $\nu$ 배선의 시계 오프셋.} 참조 적분기가 $\tau$ 를 0 에서
시작해 $\tau_0\neq0$ 인 run 에서 커널 경로와 \emph{다른 적색이동}을 보고 있었다
(상대차 $4.5\times10^{-3}$, 광학깊이 26.6 대 88.6).

\item \textbf{조용한 clamp 두 곳.} 범위 밖에서 예외 대신 clamp 하는 외부 이력이
모든 게이트를 통과하고 광학깊이를 3.3배 바꿨고, $z<0$ 요청을 눌러 $\nu$ 가
얼어붙는 경로가 있었다. 둘 다 거절로 바꿨다.

\item \textbf{파국적 상쇄 2건.} 3항 계수 $c_\lambda(x)$ 가 $x\lambda\lesssim1$ 에서,
$\Delta I$ 계수가 $x\lesssim10^{-7}$ 에서 상쇄로 무너졌다 (후자는 부호까지 틀렸다).
\texttt{expm1}/급수 형태로 교체했다.
\end{itemize}

\section{알려진 간극}

\begin{itemize}\itemsep2pt
\item \textbf{★ 원자 미시물리 (가장 큰 경계).} 이온화 이력 $x_e(z)$ 는 외부 공급이고
재결합·재이온화는 자기일관하게 풀리지 않는다.  어댑터가 이력을 \emph{소비}하는 것과
Bianchi 재결합을 \emph{푸는} 것은 완전히 다르다.  후자는 별도의 physics engine 이다.
\item \textbf{성립 영역.} Thomson 극한, 전자 정지계 산란, 단색.
Kompaneets/유도산란, 자기장(Faraday), 전자 열운동 $O(v^2)$ 은 제외.
\item \textbf{$\nu_\tau$ 의 절대 규모 --- 두 lane 으로 분리했다 (Q20).}
$H$ 앵커를 ``사용자 선택사항'' 으로 남기면 나중에 재결합 이동을 봤을 때 그것이
\emph{anisotropic recombination physics} 인지 그냥 \emph{$H$ normalization 을 다르게
준 것} 인지 섞인다.  그래서 \texttt{lane} 을 \textbf{필수 인자}로 강제했다:
\texttt{internal} 은 Gauss 구속이 정하는 $H$ 만 쓰고 외부 앵커 지정을 거절하며,
\texttt{phenomenology} 는 관측 calibrate 된 앵커를 쓰되 초기자료의 $\Omega_0$ 정합을
\textbf{예외로} 검사한다.  두 lane 은 실측 $0.389$ 배 차이 나고, 산출물에 lane 이
각인된다.
\item \textbf{$Q/U$ 다중극의 비수렴.} 스크린 기저 불연속 때문. 정량 결론에는
기저-없는 편광도만 쓴다.
\item \textbf{인접 반스텝 충돌 융합.} 비트 단위로 같은 연산자임을 확인했고 융합하면
약 25\% 빨라지지만, 상수 $\nu$ 경로의 비트 패턴이 바뀌어 골든을 재생성해야 하므로
미뤘다.
\item \textbf{$\ell=2$ 블록 부호 규약.} 내부 두 기록이 $\pm\sqrt6/10$ 로 어긋나 있다
(곱만 들어가 물리는 무해).
\end{itemize}


\section{로드맵 --- transport core 는 여기서 동결한다}

남은 문제는 더 이상 수치해석이나 기하학이 아니라 \textbf{cosmological microphysics
와 normalization} 이다.  순서를 못박는다.

\begin{enumerate}\itemsep2pt
\item \textbf{transport core 동결.} 11유형 커버리지 매트릭스를 기준선으로 고정. \emph{(완료)}
\item \textbf{$H$-anchor 문제를 먼저 닫는다} --- 원자물리보다 \textbf{앞}이다.
      \texttt{lane} 강제로 완료 (Q20).
\item \textbf{외부 정밀 재결합 provider} 연결 --- FLRW 극한에서 HyRec/Recfast 계열과
      $x_e(z)$, 가시함수, $\tau_{\rm opt}$, $z_*$ 를 맞춘다.  native atom solver 불필요.
\item \textbf{external-history + anisotropic transport lane} 을 과학적으로 완성.
\item \textbf{native Bianchi recombination} 을 별도 프로젝트로 --- multi-level atom
      + Ly$\alpha$ transfer 는 마지막 5\,\% 가 아니라 새 physics engine 하나다.
\item \textbf{재이온화} 도 같은 순서 --- 외부 $x_e(z)$/source history 먼저, native RT 나중.
\end{enumerate}

\noindent\textbf{동결 선언.} Kompaneets, Faraday, 비가환 기하 등 \emph{수송} 확장은
이 로드맵이 3--4 단계를 통과하기 전에는 착수하지 않는다.  지금 코드가 틀려서
재결합이 어려운 것이 아니라 \textbf{재결합이 원래 어려운 것} 때문에 어려운
단계로 넘어왔고, 그 경계를 흐리지 않는 것이 이 문단의 목적이다.

\section{재현}

\begin{center}\small
\begin{tabular}{@{}ll@{}}
\toprule
그림 재생성 & \texttt{python -m scripts.make\_report\_figs} \\
11유형 커버리지 & \texttt{python -m scripts.coverage\_matrix} \\
오차예산 재생성 & \texttt{python -m scripts.q17\_convergence} \\
편광 수송 유도 감사 & \texttt{python -m audit.p1\_polarized\_transport} \\
편광 Thomson 감사 & \texttt{python -m audit.p3\_polarized\_thomson} \\
부스트 유도 감사 (D7) & \texttt{python -m audit.p9\_boost\_screen} \\
회귀시험 & \texttt{pytest} 1682 passed / \texttt{cargo test --release} 122 passed \\
\bottomrule
\end{tabular}
\end{center}

\noindent 모든 임계값은 \texttt{bianchi/q/contract.py} 의 \texttt{ERROR\_BUDGET} 에서
읽는다 --- 시험이 자기 임계를 정하면 ``합격을 위한 조정''이 생기기 때문이다.
알려진 간극은 같은 파일의 \texttt{KNOWN\_GAPS} 에 기계가독으로 등재되어 있다.

\end{document}
